\Askhsh{28476}{4}
\begin{ylh}{1}
    \item 1.5 Ιδιότητες των ορίων
    \item 1.8 Συνέχεια συνάρτησης
    \item 2.1 Η έννοια της παραγώγου
    \item 2.2 Παραγωγίσιμες συναρτήσεις - Παράγωγος συνάρτηση
    \item 2.6 Συνέπειες του Θεωρήματος Μέσης Τιμής
    \item 2.7 Τοπικά ακρότατα συνάρτησης
    \item 2.9 Ασύμπτωτες - Κανόνες De l’ Hospital
    \item 3.5 Η συνάρτηση \int_a^x f(t) dt
    \item 3.7 Εμβαδόν επιπέδου χωρίου.
\end{ylh}
ΘΕΜΑ 4

Δίνεται η παραγωγίσιμη συνάρτηση $f: \mathbb{R} \to \mathbb{R}$ για την οποία ισχύουν:
\[ \lim_{x \to 1} \frac{f(x)(x-1)}{\ln x} = 0 \]
και
\[ f'(x) = \sqrt{x^2+1} \text{ για κάθε } x \in \mathbb{R}. \]
\begin{alist}
\item
    \begin{rlist}
    \item Να υπολογίσετε το $\lim_{x \to 1} \frac{\ln x}{x-1}$. \monades{03}
    \item Να αποδείξετε ότι $f(1) = 0$. \monades{03}
    \end{rlist}
\item Να αποδείξετε ότι η εξίσωση $f(x) = 0$ έχει μία ακριβώς ρίζα. \monades{06}
\item Να βρείτε το πρόσημο της συνάρτησης $f$ για κάθε $x \in \mathbb{R}$. \monades{06}
\item Να βρείτε το εμβαδόν του χωρίου $E$, που περικλείεται μεταξύ της γραφικής παράστασης της συνάρτησης $f$, τον άξονα $x'x$ και των ευθειών $x = 0$ και $x = 1$. \monades{07}
\end{alist}